\section{Model and estimation}
\label{sec:methods}

\paragraph{Likelihood model.} Let $\phi$ be a frozen CLIP ViT-B/16
encoder \citep{radford2021clip} and $g_\theta : \mathbb{R}^{512} \to
\mathcal{M}$ a two-layer MLP with GELU and dropout, where the latent
manifold $\mathcal{M}$ is either Euclidean $\mathbb{R}^d$ or a
Poincar\'e ball $\mathbb{B}^d_c$ of curvature $c > 0$. For an image
$x$ of style $y \in \{1, \ldots, K\}$, set $z = g_\theta(\phi(x))$,
and let $\{\mathbf{p}_k\}_{k=1}^K \subset \mathcal{M}$ be learnable
class prototypes. The class-conditional likelihood is a softmax over
distance to prototype on $\mathcal{M}$:
\[
  p(y = k \mid z;\, \{\mathbf{p}_k\})
    \;\propto\; \exp\!\bigl(-d_\mathcal{M}(z, \mathbf{p}_k)\bigr).
\]
The encoder feeds the head identically in both geometries; switching
geometries changes exactly two things: the head's final transform
(identity for $\mathbb{R}^d$; $\exp_0^c(u) = \tanh(\sqrt{c}\,\|u\|)\,
u/(\sqrt{c}\,\|u\|)$ for the ball) and the metric used by the
classifier:
\[
  d^c(x, y) = \tfrac{1}{\sqrt{c}}\operatorname{arcosh}\!\Bigl(
    1 + \tfrac{2c\,\|x-y\|^2}{(1-c\|x\|^2)(1-c\|y\|^2)}
  \Bigr).
\]
Backbone, MLP width, dropout, batch size, optimiser, and schedule
are shared.

\paragraph{Tree-structured regularizer.} We add a tree-structured
regularizer on prototype layouts. Let
$D_{ij} = d_\mathcal{M}(\mathbf{p}_i, \mathbf{p}_j)$ be the matrix of
pairwise prototype distances on $\mathcal{M}$ and
$T_{ij} = d_T(i, j)$ the tree distance matrix on the reference tree
$T$. The regularizer penalises deviation between the \emph{shapes} of
these two distance matrices,
\[
  \mathcal{R}(\{\mathbf{p}_k\}) \;=\;
    \frac{\lambda}{|\mathcal{P}|}
    \sum_{(i,j) \in \mathcal{P}}
    \Bigl(\tfrac{D_{ij}}{\bar{D}} - \tfrac{T_{ij}}{\bar{T}}\Bigr)^{\!2},
\]
where $\mathcal{P}$ is the set of $\binom{K}{2}$ off-diagonal pairs,
$\bar{D}, \bar{T}$ are the means of $D, T$, and $\lambda$ controls
regularizer strength. Mean-normalising each pair before differencing
makes the regularizer invariant to the absolute scale of distances on
$\mathcal{M}$: the geometry is free to pick whatever scale the
likelihood prefers, while the regularizer constrains only the
relative structure of the prototype distance matrix. $\lambda = 0$
recovers the cross-entropy baseline used in prior hyperbolic-image
work \citep{khrulkov2020hyperbolic}. We treat $\mathcal{R}$ as a
penalty rather than a Bayesian log-prior: $\mathcal{R}$ is not the
log of a normalised density on $\mathcal{M}^K$ (it is mean-rescaled
on each evaluation), so it has no proper-prior interpretation, and we
report regularized maximum-likelihood estimates rather than MAP
estimates of a posterior.

\paragraph{Estimation.} We minimise the regularized negative
log-likelihood
$\mathcal{L}(\theta, \{\mathbf{p}_k\}) = \mathcal{L}_{CE}(\theta, \{\mathbf{p}_k\}) + \mathcal{R}(\{\mathbf{p}_k\})$
by stochastic gradient descent. Adam optimises the MLP head;
Riemannian Adam from \texttt{geoopt} \citep{kochurov2020geoopt}
optimises the hyperbolic prototypes, projecting each update back onto
the manifold. The regularizer is computed once per gradient step over
the $\binom{27}{2} = 351$ off-diagonal style pairs, negligible cost
next to the per-batch likelihood term. The framing makes explicit
that the regularizer is a structured penalty on a geometric quantity
(the prototype distance matrix) and that the choice of latent
manifold controls how cheaply that penalty can be driven to zero.

\paragraph{Distortion bound.} The regularizer is a
function of pairwise prototype distances, so in any latent geometry
where the tree distance matrix $T$ admits a low-distortion isometric
embedding, the regularizer cost at the optimum is small. Sala
et~al.\ give distortion bounds $\sim 1/d$ for $\mathbb{B}^d_c$ versus
a constant floor for $\mathbb{R}^d$ at fixed $d$
\citep{sala2018representation}. The empirical question is whether
that asymptotic statement is visible on a real dataset at modest
$d$, against a real classification likelihood, with a finite-sample
fit rather than an isometric embedding objective.

\paragraph{Training details.} 150 runs total; batch size 4096; Adam
$\eta = 10^{-3}$ for the head and $10^{-2}$ for the prototypes;
weight decay $10^{-4}$; dropout $0.1$; gradient clipping at norm
$1.0$; 30 epochs. CLIP features are pre-cached as float16 tensors,
keeping each run under 10 seconds on a single Apple M-series GPU and
the full sweep under half an hour. Every headline configuration is
replicated across three seeds; reported error bars are seed standard
deviation. Significance comparisons across seeds use paired-$t$ tests
on the 18 (dim, seed) pairs of the dimension sweep, with hyperbolic
configurations selected at the best curvature per pair.

\paragraph{Evaluation.} A useful style embedding can be useful along
three orthogonal axes, and we report metrics for each.
\emph{Classification} is top-1, top-5, and balanced accuracy.
\emph{Global tree fidelity} is the Spearman correlation between the
learned prototype distance matrix and the tree distance matrix, plus
mean and worst-case multiplicative distortion. \emph{Local tree
preservation} is sibling and cousin recall@$k$ among the $k$ nearest
neighbours of each validation embedding, for $k \in \{5, 10\}$.
Sibling/cousin sets are derived from a reference tree; we evaluate
against the default, the CLIP-empirical, and the DINOv2-empirical
tree to test sensitivity (\cref{ssec:local,ssec:robust}).
